\documentclass{practice}

\title{7}
\date{\today}

%\usepackage{etoolbox}
%\makeatletter
%\preto{\@verbatim}{\topsep=0pt \partopsep=-2\parsep }
%\makeatother

\usepackage{crysymb}

\usepackage{fancyvrb}
\fvset{listparameters=\setlength{\topsep}{0pt}\setlength{\partopsep}{0pt}}

\usepackage{listings}

\lstset{
basicstyle=\small\ttfamily,
columns=flexible,
breaklines=true
}

\begin{document}
\maketitle

\begin{task}{ASN.1 DER}
  Install \texttt{dumpasn1}, generate an ECC or RSA keypair using OpenSSL, and visualise what the key files really contain.
  Note that you must strip the PEM header and decode the base64 data before \texttt{dumpasn1} can parse the ASN.1 bytes.

  Now use OpenSSL's \href{https://docs.openssl.org/master/man1/openssl-asn1parse/}{\texttt{asn1parse}} tool to do the same thing.
  Which output do you prefer?
\end{task}

\begin{task}{ElGamal}
  \textit{Preface.}
  In the lecture we saw that ElGamal is \emph{partially homomorphic} with respect to multiplication, often simply called \enquote{multiplicatively homomorphic}.
  That is, if you multiply two ciphertexts $c_1$, $c_2$ and decrypt the resulting ciphertext, you will get $m_1 \cdot m_2$.
  More mathematically,
  \begin{align*}
    \DEC_\SK(c_1 \cdot c_2) &= 
    \DEC_\SK\bigl(\ENC_\PK(m_1) \cdot \ENC_\PK(m_2)\bigr)\\
    &= \DEC_\SK\bigl(\ENC_\PK(m_1\cdot m_2)\bigr)\\
    &= m_1 \cdot m_2
  \end{align*}
  holds for classic ElGamal.

  In additive notation, e.g. as for elliptic curves, multiplication becomes addition and exponentiation becomes multiplication.
  By convention, we denote group elements (points on a curve) by capital letters and scalars (exponents in multiplicative notation) by lowercase letters.

  The homomorphism then becomes:
  \begin{align*}
    \DEC_\SK(C_1 + C_2) &= 
    \DEC_\SK\bigl(\ENC_\PK(M_1) + \ENC_\PK(M_2)\bigr)\\
    &= \DEC_\SK\bigl(\ENC_\PK(M_1 + M_2)\bigr)\\
    &= M_1 + M_2
  \end{align*}

  We now face an issue: how do we convert a message (a bytestring) into a point on the elliptic curve?
  While there are techniques for this, let us use lifted ElGamal instead.
  In lifted ElGamal, the message is the scalar (an integer), and the message bytes can then simply be interpreted as an integer.
  The EC lifted ElGamal homomorphism is:
  \begin{align*}
    \DEC_\SK(C_1 + C_2) &= 
    \DEC_\SK\bigl(\ENC_\PK(m_1G) + \ENC_\PK(m_2G)\bigr)\\
    &= \DEC_\SK\bigl(\ENC_\PK((m_1 + m_2)G)\bigr)\\
    &= (m_1 + m_2)G
  \end{align*}

  \textit{Environment.}
  I like to use \href{https://pypi.org/project/fastecdsa/#arbitrary-elliptic-curve-arithmetic}{\texttt{fastecdsa}} to experiment with elliptic curve operations in Python.
  Before you can install it on Linux or macOS, make sure that you have installed \href{https://gmplib.org}{\textit{GMP}}, e.g. with
  \begin{itemize}
    \item \texttt{apt install libgmp3-dev}
    \item \texttt{brew install gmp}
  \end{itemize}
  It is not well suited for Windows, but it should work in WSL.
  
  On macOS you may need to install \href{https://github.com/AntonKueltz/fastecdsa/issues/74#issuecomment-932757298}{\texttt{fastecdsa}} with
  \begin{Verbatim}
CFLAGS=-I/opt/homebrew/opt/gmp/include \
  LDFLAGS=-L/opt/homebrew/opt/gmp/lib \
  pip install fastecdsa
  \end{Verbatim}

  Alternatively, for less \enquote{hassle}, you can use \href{https://pypi.org/project/tinyec/}{\texttt{tinyec}}.

  \textit{Task.}
  For this task, generate a P-384 keypair and encrypt two different integers using lifted EC ElGamal.
  Add the ciphertexts and decrypt the result.
  Is the result the sum of the two integers?

  Then, repeat three times: \enquote{my implementation is insecure and I will not use it for actual encryption}.
\end{task}

\begin{task}{Sage RSA}
  SageMath\footnote{\url{https://www.sagemath.org}} is a computer algebra system (CAS) which can be programmed in with a Python-based language.
  It is free and open-source as opposed to Wolfram Mathematica and MATLAB, and is very powerful.
  As such, it is quite popular among cryptographers for some quick computations/prototyping, and is also popular for cryptography tasks in CTFs (Capture The Flag exercises).

  Check out also Sage's \href{https://doc.sagemath.org/html/en/reference/cryptography/index.html}{\enquote{cryptography reference}}.

  \textit{Task.}
  Given the following RSA public key:
  \begin{itemize}
    \item $e = 65537$
    \item $n = 230624063920680065513680711198246109591$,
  \end{itemize}
  decipher the RSA ciphertext $c = 161670843286605628968182075437059446607$.

  Note that the RSA modulus is $128$ bits long, so \enquote{naively} brute forcing either prime is not a good approach.
  Can you think of a more efficient approach?
  
  \textit{Hint 1.}
  RSA-100 ($330$ bits) was \emph{factored} way back in 1991.

  \textit{Hint 2.}
  Recovering the message takes less than a second on my laptop with the proper Sage tools.

  \textit{Addendum.}
  Check out the RsaCtfTool\footnote{\url{https://github.com/RsaCtfTool/RsaCtfTool}}.
  Can you solve the task using the tool?
\end{task}
\end{document}
